% A.72 candidate: notes only, awaiting full independent audits.
\subsection{Every generalized-quaternion geometric closure}

\begin{statement}{A.72}{Proved}{Both cusp assertions for every generalized-quaternion closure}
Let \(Y\to U=E\setminus\{O\}\) be an actual connected finite
\'{e}tale cover, where \(E:y^2=x^3-x\), whose geometric Galois
closure has group
\[
 Q_{4n}=\langle r,s\mid r^{2n}=1,\ s^2=r^n,
                         \ srs^{-1}=r^{-1}\rangle,
 \qquad n\ge2\text{ a power of two}.                  \tag{A.72.1}
\]
Then both assertions of Task A.5 hold for \(Y\), at every
coefficient prime, every induced cusp packet, and every permitted
open derived subgroup. In particular the full fixed geometric
pro-\(\ell\) subgroup is the selected cusp inertia, and every
nontrivial arithmetically normalized procyclic pro-\(\ell\)
subgroup is cuspidal.

More precisely, the smooth proper Galois closure \(C\) has genus
\(2n-1\), four cusps, potential good Jacobian reduction of
\(2\)-rank zero at two, and every cusp difference is killed by
\(2n\). Every actual intermediate inherits the good-reduction
and torsion-divisor properties needed for the fixed-group
conclusion. This settles this family of geometric closures;
the universal question A.45 remains open.
\end{statement}

\begin{proof}
We start with the actual geometric Galois closure of the given
cover. After a finite number-field extension it and its deck
transformations, quotient maps, and cusps are defined over that
field. No identification with an unproved explicit curve model
is needed. All further reduction arguments are at a place over
two, after a finite extension making the elliptic base good,
the necessary stable models available, and all endomorphisms
defined. Local automorphisms fix the coefficient valuation ring.

\paragraph{The actual inertia subgroup and the four cusps.}
Write elements of \(G=Q_{4n}\) uniquely as \(r^a s^b\), with
\(a\in\mathbf Z/(2n)\) and \(b\in\{0,1\}\). Their product is
\[
 (a,b)(c,d)=\bigl(a+(-1)^b c+n b d,\ b+d\bmod2\bigr),
                                                        \tag{A.72.2}
\]
where the first coordinate is reduced modulo \(2n\).
The geometric fundamental group of the once-punctured elliptic
curve is free profinite on two generators; its peripheral
generator is their commutator, up to inverse and conjugacy.
Thus the image of a generating pair in this actual quotient
determines its actual cusp inertia.

A generating pair cannot consist of two rotations. If one entry
is \(r^a\) and the other \(r^b s\), it generates precisely
when \(a\) is odd. If both entries are outside
\(\langle r\rangle\), say \(r^a s,r^b s\), it generates
precisely when \(a-b\) is odd: their quotient is a rotation
with that exponent, and their squares are \(r^n\).
These assertions follow because \(2n\) is a power of two.
For the commutator convention \([x,y]=xyx^{-1}y^{-1}\),
\[
 [r^a,r^b s]=r^{2a},\qquad
 [r^a s,r^b s]=r^{2(a-b)}.                         \tag{A.72.3}
\]
Reversing the entries in the first formula takes the inverse.
Consequently every generating pair has commutator generating
the same normal subgroup
\[
                 I=\langle r^2\rangle,
 \qquad |I|=n,\qquad G/I\simeq C_2^2.             \tag{A.72.4}
\]
This is therefore the actual inertia at every cusp. There are
four cusps and their ramification indices are all \(n\).
Riemann--Hurwitz for the actual degree-\(4n\) map \(C\to E\)
gives
\[
                    2g(C)-2=4(n-1),
                    \qquad g(C)=2n-1.             \tag{A.72.5}
\]

The same elementary description proves the needed Frattini
property. The squares and commutators generate \(I\), since
all outside elements square to \(r^n\), while the square of
\(r\) is \(r^2\). Equivalently, if a subgroup \(H\) satisfies
\(HI=G\), it contains an odd rotation and an outside element,
which generate \(G\). Hence \(I=\Phi(G)\), and no proper
subgroup surjects onto \(G/I\). This includes \(n=2\), where
the group is \(Q_8\).

\paragraph{Compact type and the actual central quotient.}
The group \(G\) is nilpotent. The elliptic base has potential
good reduction at two, the sole branch section \(O\) is
\'{e}tale over the valuation ring, and generic ramification is
tame in characteristic zero. Lehr--Matignon, Theorem 6.13,
therefore gives potential good reduction of \(J_C\), or
equivalently a tree as the dual graph of the proper stable
fibre. We use exactly this ramified nilpotent-cover conclusion,
also recorded by Sa\"idi on printed page 662.
\cite{fghLehrMatignon2005,fghSaidi2007}
No slope statement for an exceptional component over a marked
branch point is inferred from it.

The unique involution of \(G\) is
\(\sigma=r^n\). It is central and lies in \(I\).
The actual quotient \(C_0=C/\langle\sigma\rangle\) is a
cover of the same fixed elliptic base, with group
\(G/\langle\sigma\rangle\simeq D_{2n}\), inertia of order
\(n/2\), four cusps, and genus \(n-1\). A.70 identifies
every such actual dihedral cover with one of its three models;
the reduction assertions of A.44 and A.69 show that
\(J_{C_0}\) has potential good reduction with special
\(2\)-rank zero. When \(n=2\), the quotient group is
\(V_4\), and \(C_0\to E\) is an elliptic isogeny, giving
the same conclusion from the supersingularity of the good
reduction of \(E\) at two.

\paragraph{The anti-invariant factor has only two possible ranks.}
Let \(A\) be the anti-invariant isogeny factor for \(\sigma\).
The actual pullback and norm maps give
\[
             J_C\sim J_{C_0}\times A,
             \qquad \dim A=n.                    \tag{A.72.6}
\]
On \(A\) we have \(r^n=-1\). Since \(n\) is a power of
two, \(x^n+1=\Phi_{2n}(x)\) is irreducible over \(\mathbf Q\).
The action is unital and \(A\ne0\), so it gives an injection
\[
             \mathbf Q(\zeta_{2n})
                       \hookrightarrow\operatorname{End}^0(A).
                                                        \tag{A.72.7}
\]
This argument uses the actual deck endomorphism, without any
assumption about a Hodge character decomposition.

After the indicated good-reduction extension, endomorphisms
extend to the abelian scheme. Thus the field
\(\mathbf Q_2(\zeta_{2n})\), of degree \(n\) over
\(\mathbf Q_2\), acts with its unit on the rational \'{e}tale
Tate space of the special \(2\)-divisible group of \(A\).
If this space is nonzero its dimension is divisible by \(n\).
Its dimension is the special \(2\)-rank \(f(A)\), at most
\(\dim A=n\). Hence
\[
                    f(A)=0\quad\hbox{or}\quad f(A)=n.
                                                        \tag{A.72.8}
\]
The second alternative means that \(A\) is ordinary. We now
exclude it by the actual action on stable components.

\paragraph{The component calculation under the ordinary alternative.}
Suppose \(A\) is ordinary, and choose the stable marked model
of \(C\) after the finite extension already described. Over
an algebraically closed residue field its graph is a tree, and
its good special Jacobian is the product of the Jacobians of
its positive-genus components. The specialization of the
\(\sigma\)-invariant factor has \(2\)-rank zero, whereas
the anti-invariant factor is ordinary.

If \(\sigma\) exchanges two positive-genus components, their
invariant diagonal and anti-invariant diagonal are both
isogenous to the same component Jacobian. A positive-dimensional
abelian variety cannot have both rank zero and be ordinary,
so no such pair occurs. A component on which \(\sigma\)
acts identically contributes only to the invariant factor;
we do not infer that its genus is zero.

Let \(B\) be a component fixed by \(\sigma\), with a
nontrivial induced involution, and put \(B'=B/\sigma\).
The actual pullback--norm identities identify the invariant
Jacobian factor with \(J_{B'}\) up to isogeny, hence
\(f(B')=0\). Its anti-invariant Prym is ordinary when
positive-dimensional. Write \(r_B\) for the number of fixed
points of \(\sigma\) on \(B\), and \(d_P\) for their
different exponents. The smooth-curve Deuring--Shafarevich
and Riemann--Hurwitz formulas give
\[
 \begin{aligned}
 f(B)&=r_B-1,\\
 g(B)-g(B')&=g(B')-1+\frac12\sum_{P\ {\rm fixed}}d_P,
                    \qquad d_P\ge2.
 \end{aligned}                                         \tag{A.72.9}
\]
Here \(r_B\ge1\), since otherwise the first equation gives
negative rank. If the Prym has dimension zero, these equations
force \(g(B')=g(B)=0\) and \(r_B=1\). If it has positive
dimension, its ordinarity implies
\(f(B)=g(B)-g(B')\). Therefore
\[
 r_B=g(B')+\frac12\sum d_P
       \quad\Longrightarrow\quad
 g(B')=0,\qquad d_P=2\text{ for every fixed point}.
                                                        \tag{A.72.10}
\]
Every positive-genus component contributing to the anti-invariant
factor is thus ordinary; its whole Jacobian is anti-invariant,
and \(\sigma\) has genus-zero quotient. The sum of the genera
of these components is exactly \(\dim A=n\).
The cited Deuring--Shafarevich formula applies to the smooth
component curves, not to an arbitrary marked vertical fibre.
\cite[Corollary 1.8]{fghCrew1984}
\cite[Proposition 1.4]{fghYang2024}

\paragraph{Why a fourth root cannot act on such a component.}
We repeat the elementary ramification argument at its required
scope. For a faithful finite \(2\)-group \(T\) acting on an
ordinary smooth proper curve \(B\), subtract twice
Deuring--Shafarevich from Riemann--Hurwitz. With branch indices
\(e_Q\) and different exponents \(d_Q\), it gives
\[
 0=2|T|\bigl(g(B/T)-f(B/T)\bigr)
      +\sum_Q\frac{|T|}{e_Q}
                         \bigl(d_Q-2(e_Q-1)\bigr).
                                                        \tag{A.72.11}
\]
All summands are nonnegative. Indeed, for wild \(2\)-group
inertia the classical lower groups have \(T_0=T_1\), and
\(d_Q=\sum_{j\ge0}(|T_j|-1)\ge2(e_Q-1)\).
The equality forces \(T_2=1\) at each branch point. The
coefficient of \(t^2\) in the local action
\(t\mapsto t+c t^2+\cdots\) embeds \(T_1/T_2\) in the
additive residue field. Thus a cyclic subgroup of order four
cannot fix a point on an ordinary curve. This uses the
classical equal-characteristic lower numbering.
\cite[\S2.2, equation (3) and Exercise 2.12(c)]{fghBouwWewers}

If an order-four automorphism \(a\) has square \(\sigma\)
and \(B/\sigma\simeq\mathbf P^1\), its induced involution
on \(\mathbf P^1\) has a unique fixed point \(Q\) in
characteristic two. This point cannot be unramified in
\(B\to B/\sigma\): on its two-point fibre the square of
any permutation is trivial, whereas \(a^2=\sigma\) would
exchange the points. Hence the fibre is a single ramified
point fixed by \(a\), contradicting (A.72.11). There is
therefore no such fourth root on an ordinary component.

\paragraph{The unique involution gives the contradiction.}
Take a positive-genus anti-invariant component \(B\), which
exists because \(\dim A=n>0\). Its stabilizer \(G_B\)
contains \(\sigma\). Its action on \(B\) is faithful:
every nontrivial subgroup of the quaternion \(2\)-group
contains its unique involution \(\sigma\), whereas
\(\sigma\) acts nontrivially on \(B\).
For completeness, a nonidentity rotation has a power equal
to \(r^n\), and every outside element has square \(r^n\);
this proves the subgroup assertion directly.

If \(G_B\) were larger than \(\langle\sigma\rangle\),
it would contain an element of order at least four. A power
of that element has order four and square \(\sigma\).
Faithfulness would give precisely the fourth root just
excluded. Thus \(G_B=\langle\sigma\rangle\), and the
orbit of \(B\) has \(2n\) components. Since \(\sigma\)
is central, all these components are anti-invariant and
have positive genus. Their genera already sum to at least
\(2n\), contrary to the total anti-invariant dimension \(n\).
This excludes the ordinary alternative in (A.72.8).
Consequently \(f(A)=0\), and \(J_C\) has special
\(2\)-rank zero. No smooth fixed-point lifting assertion
or unproved marked-fibre slope assertion enters this argument.

\paragraph{The actual elliptic quotient gives cusp torsion.}
Let \(E'=C/I\), with its actual map to \(E\). All cusp
inertia is killed in this quotient, so it extends to a
connected unramified degree-four cover of the proper
elliptic curve. Choose a point over \(O\) as its origin.
It is then an elliptic isogeny with kernel \(C_2^2\), hence
with kernel exactly \(E'[2]\). Write its four kernel points
as \(Q_1,\ldots,Q_4\), and the cusps of \(C\) above them
as \(P_1,\ldots,P_4\).

The actual map \(h:C\to E'\) has degree \(n\). At each
\(Q_j\) its inertia is the whole group \(I\), so there is
one point above it and
\[
                     h^*(Q_j)=nP_j.
                                                        \tag{A.72.12}
\]
Every \(Q_i-Q_j\) is a two-torsion divisor class on \(E'\).
Choose an actual rational function \(f_{ij}\) with
\(\operatorname{div}(f_{ij})=2(Q_i-Q_j)\), after the same
permitted finite extension. Pulling it back gives
\[
       \operatorname{div}(f_{ij}\circ h)=2n(P_i-P_j).
                                                        \tag{A.72.13}
\]
Thus the required torsion comes from the actual intermediate
elliptic map, with its total ramification, not just from an
isogeny factor of the Jacobian.

\paragraph{Every actual intermediate and the arithmetic conclusion.}
Let \(H<G\) be proper and \(D=C/H\). Because \(I\) is
normal, its cusps are the \(H\)-orbits on \(G/I\), so
their number is \([G:HI]\ge2\). Its compact genus is
positive, since it has a nonconstant finite map to \(E\).
The actual norm map makes \(J_D\) an isogeny quotient
of \(J_C\), preserving potential good reduction and special
\(2\)-rank zero. For any two cusps of \(D\), choose cusps
above them in \(C\) and take the norm of the function in
(A.72.13). Over an algebraically closed field the pushforward
of a point is its image with coefficient one, regardless
of ramification. Hence every target cusp difference is
also killed by \(2n\).

In the compact-type marked model of \(D\), these cusp
differences specialize to zero, because an abelian variety
of \(2\)-rank zero has no nonzero geometric \(2\)-power
torsion points. The extended Abel map and A.41 therefore
give a rational joining path from every selected cusp to
any distinct companion cusp. A.42 proves the complete
fixed-group assertion for every coefficient prime, including
two. A.59 gives all the normalizer assertions, since the
actual geometric Galois closure is a two-group. The quotient
by \(H=G\) is the already settled original curve \(U\).

Finally finite field extensions are used only to define
curves, deck transformations, subgroup quotients, and marked
points and to obtain reduction. A geometric path identifying
base points need not be defined over a finite extension;
the resulting comparison is up to geometric inner conjugacy,
which preserves the packet and inertia formulations. The
restriction argument in A.41--A.42 handles the original
permitted open derived subgroup, and normalizers descend by
intersecting their open arithmetic image with the field
extension. At no point is fixedness transferred through a
noninjective geometric quotient.

The exact multiplication, commutator, subgroup, cusp-count,
cyclotomic-dimension and component-orbit checks are retained
in \path{notes/fgh_A72_certificate.py}. The written arguments
establish the assertion for every power of two, independently
of the finite verification range.
\end{proof}
